\documentclass[a4paper,6pt,twoside]{jsarticle}
\usepackage[utf8]{inputenc}
\usepackage{amsmath, amssymb}
\usepackage{enumitem}
\usepackage[top=10mm,bottom=10mm,left=10mm,right=10mm]{geometry}
\pagestyle{empty}

% 〇×の記号表示マクロを定義
\newcommand{\OX}{\hfill（〇／×）}

% 行間とスペースを詰める設定
\setlength{\baselineskip}{8pt}
\setlength{\parskip}{0pt}
\setlength{\parindent}{0pt}
\setlist[enumerate]{itemsep=0pt, parsep=0pt, topsep=0pt, leftmargin=15pt}
\setlist[itemize]{itemsep=0pt, parsep=0pt, topsep=0pt, leftmargin=15pt}

\begin{document}

\begin{center}
\normalsize{情報リテラシー 2025年度前期末試験 問題用紙} \quad 出席番号：\underline{\hspace{1.5cm}} \quad 氏名：\underline{\hspace{6cm}}
\end{center}

\vspace{0mm}

\subsubsection*{注意事項}
{\small
\begin{itemize}
  \item マークシートの学年は、1を記入・マークしてください。
  \item マークシートのクラスは、M:01、E:02、I:03、C:04、A:05の該当する数字を記入・マークしてください。
  \item マークシートの番号は、出席番号を3桁に変換して記入・マークしてください。
  \item 問題は全部で \textbf{50問} あります。各問題は \textbf{2点} です。\textbf{100点満点} です。
  \item 解答は\textbf{解答用紙}の解答欄に \textbf{正確に記入}してください。枠外の解答は採点対象にはなりません。
  \item 問1〜問4は〇×問題です。内容が正しい場合は〇、誤っている場合は×を選び、解答用紙に記入すること。
  \item 問5〜問7は三択問題です。各設問について最も適切な選択肢を1つ選び、A・B・Cのいずれかを解答用紙に記入すること。
  \item 到達目標1：問1, 問3, 問4, 問5, 問6, 問7、到達目標3：問2。
\end{itemize}
}

\vspace{1mm}

\subsubsection*{問1：情報の収集・整理１（各2点×7問＝14点）}
\begin{enumerate}
\item メールの件名には用件と名前を入れることが推奨される。\OX
\item メールはチャットよりもフォーマルな文体で書くべきである。\OX
\item メールには署名を付けることが推奨される。\OX
\item 緊急の連絡にはチャットが適している。\OX
\item CCとBCCは同じ機能である。\OX
\item メールは「一文一情報」で書くことが大切である。\OX
\item チャットとメールでは、メールの方が記録として残りやすい。\OX
\end{enumerate}

\subsubsection*{問2：情報の収集・整理2（各2点×7問＝14点）}
\begin{enumerate}
\item インターネット上のすべての情報は正確である。\OX
\item 情報の真偽を見分けるには、どこから得た情報かを確認することが大切である。\OX
\item フェイクニュースは現代社会に存在しない。\OX
\item 情報の正確性を判断するには複数の出典を確認すると良い。\OX
\item SNSで見た情報は必ず本当である。\OX
\item デマ情報を拡散することには責任が伴う。\OX
\item 自分に都合の良い情報だけを信じるのは危険である。\OX
\end{enumerate}

\subsubsection*{問3：情報の加工１（各2点×7問＝14点）}
\begin{enumerate}
\item PowerPointファイルの拡張子は「.pptx」である。\OX
\item ファイルの拡張子「.docx」はWord文書の形式である。\OX
\item Wordのファイルは自動保存機能が使える。\OX
\item Wordでは文字の書式設定（フォント、サイズ、色など）を変更できる。\OX
\item Wordで作成した文書はPDF形式で保存することができる。\OX
\item 学校のネットワーク負荷を軽減するため、基本的にはローカル保存を行うべきである。\OX
\item ローカル保存はクラウド保存よりも高速でアクセスできる。\OX
\end{enumerate}

\subsubsection*{問4：コンピュータのしくみ1（各2点×7問＝14点）}
\begin{enumerate}
\item コンピュータは入力・処理・出力・記憶・制御の5つの装置から成る。\OX
\item メモリは一時的にデータを記憶する装置である。\OX
\item ハードディスクは処理装置である。\OX
\item 制御装置は他の装置を指示・制御する。\OX
\item 入力装置は出力結果を表示する。\OX
\item CPUは中央処理装置とも呼ばれ、演算や制御を行う。\OX
\item キーボードやマウスは入力装置の例である。\OX
\end{enumerate}

\subsubsection*{問5：コンピュータのしくみ2（各2点×7問＝14点）}
\begin{enumerate}
\item OSの主な役割として最も適切なのはどれか。\\
A.~文書を作成する \quad B.~ハードウェアとアプリケーションを繋ぐ \quad C.~インターネットに接続する

\item システムソフトウェアに分類されるのはどれか。\\
A.~Word \quad B.~Chrome \quad C.~OS

\item CLIの特徴として最も適切なのはどれか。\\
A.~マウスで操作する \quad B.~文字でコンピュータに指示する \quad C.~画面をタッチして操作する

\item スマートフォンの主なインターフェースはどれか。\\
A.~CLI \quad B.~GUI \quad C.~タッチインターフェース

\item 有線接続の特徴として最も適切なのはどれか。\\
A.~配線が不要 \quad B.~接続が安定している \quad C.~電力消費が少ない

\item Bluetoothの特徴として最も適切なのはどれか。\\
A.~高品質な映像出力ができる \quad B.~近距離の機器接続で低電力 \quad C.~インターネット接続に使用する

\item アプリケーションソフトウェアの例として正しいのはどれか。\\
A.~Windows \quad B.~デバイスドライバ \quad C.~Excel
\end{enumerate}

\vspace{-3mm}

\subsubsection*{問6：情報のデジタル表現1（各2点×7問＝14点）}
\begin{enumerate}
\item アナログとデジタルの違いとして最も適切なのはどれか。\\
A.~アナログは連続的、デジタルは離散的 \quad B.~アナログは新しい、デジタルは古い \quad C.~アナログは速い、デジタルは遅い

\item 2進数「101」を10進数に変換すると何になるか。\\
A.~3 \quad B.~5 \quad C.~7

\item 1バイトは何ビットか。\\
A.~4ビット \quad B.~8ビット \quad C.~16ビット

\item 論理積（AND）演算で、1 AND 0の結果はどれか。\\
A.~0 \quad B.~1 \quad C.~2

\item 文字コードの説明として最も適切なのはどれか。\\
A.~文字を数値で表現するための対応表 \quad B.~文字を画像に変換する方法 \quad C.~文字を音声に変換する技術

\item 16進数「A」を10進数に変換すると何になるか。\\
A.~9 \quad B.~10 \quad C.~11

\item 1MBは何バイトか。\\
A.~1,000,000バイト \quad B.~1,024,000バイト \quad C.~1,048,576バイト
\end{enumerate}

\subsubsection*{問7：情報のデジタル表現2（各2点×8問＝16点）}
\begin{enumerate}
\item 4桁の2進数で「1010-0011」を2の補数を使って計算した結果はどれか。\\
A.~0111 \quad B.~1000 \quad C.~1111

\item 2の補数を使う理由として最も適切なのはどれか。\\
A.~足し算だけで引き算ができる \quad B.~計算が速くなる \quad C.~メモリを節約できる

\item 浮動小数点演算で生じる問題はどれか。\\
A.~計算が遅くなる \quad B.~丸め誤差が発生する \quad C.~メモリを多く使う

\item 音声のデジタル化の手順として正しいのはどれか。\\
A.~符号化→量子化→サンプリング \quad B.~サンプリング→量子化→符号化 \quad C.~量子化→符号化→サンプリング

\item CD音質の規格として正しいのはどれか。\\
A.~8kHz、8ビット \quad B.~44.1kHz、16ビット \quad C.~96kHz、24ビット

\item RGBで白色を表現する値はどれか。\\
A.~R=0, G=0, B=0 \quad B.~R=255, G=255, B=255 \quad C.~R=128, G=128, B=128

\item 可逆圧縮の特徴として最も適切なのはどれか。\\
A.~ファイルサイズが大幅に小さくなる \quad B.~元のデータを完全に復元できる \quad C.~処理速度が速い

\item JPEGファイルの特徴として最も適切なのはどれか。\\
A.~無圧縮で高画質 \quad B.~可逆圧縮で中サイズ \quad C.~非可逆圧縮で小サイズ
\end{enumerate}

\end{document}